\documentclass[11pt,a4paper]{article}

\usepackage[utf8]{inputenc}
\usepackage[T1]{fontenc}
\usepackage{lmodern}
\usepackage[margin=2.3cm]{geometry}
\usepackage[table]{xcolor}
\usepackage{listings}
\usepackage{enumitem}
\usepackage{titlesec}
\usepackage{parskip}
\usepackage[hidelinks]{hyperref}
\usepackage{fancyhdr}
\usepackage{tabularx}
\usepackage{array}

% ---------- Colours (matching the setup manual) ----------
\definecolor{codebg}{HTML}{F4F4F4}
\definecolor{codeborder}{HTML}{DDDDDD}
\definecolor{codecomment}{HTML}{6A737D}
\definecolor{codekey}{HTML}{0550AE}
\definecolor{codestr}{HTML}{0A7F3F}
\definecolor{accent}{HTML}{0F4F8B}
\definecolor{warnbg}{HTML}{FFF4E5}
\definecolor{warnbd}{HTML}{D58512}
\definecolor{notebg}{HTML}{EAF4FB}
\definecolor{notebd}{HTML}{2E7CB7}

\lstset{
  backgroundcolor=\color{codebg},
  basicstyle=\ttfamily\footnotesize,
  breaklines=true,
  breakatwhitespace=false,
  columns=fullflexible,
  frame=single,
  rulecolor=\color{codeborder},
  framesep=4pt,
  xleftmargin=4pt,
  xrightmargin=4pt,
  showstringspaces=false,
  commentstyle=\color{codecomment},
  keywordstyle=\color{codekey}\bfseries,
  stringstyle=\color{codestr},
  literate={"}{{\textquotedbl}}1
}

\lstdefinelanguage{js}{
  morekeywords={const,let,var,function,return,await,async,if,else,new,from,import,export,try,catch,throw},
  morecomment=[l]{//},
  morecomment=[s]{/*}{*/},
  morestring=[b]",
  morestring=[b]'
}

\titleformat{\section}{\Large\bfseries\color{accent}}{\thesection}{0.6em}{}
\titleformat{\subsection}{\large\bfseries\color{accent!85!black}}{\thesubsection}{0.5em}{}
\titlespacing*{\section}{0pt}{1.6em}{0.6em}
\titlespacing*{\subsection}{0pt}{1.2em}{0.4em}

\pagestyle{fancy}
\fancyhf{}
\renewcommand{\headrulewidth}{0pt}
\fancyfoot[C]{\thepage}
\fancyfoot[L]{\small\color{gray}Android stale-cache diagnosis --- 2026-06-11}

\newcommand{\callout}[3]{%
  \begin{center}%
  \begin{tabularx}{\linewidth}{|>{\columncolor{#1}}X|}%
  \hline
  \textbf{#2}\par\smallskip #3\\[2pt]
  \hline
  \end{tabularx}%
  \end{center}%
}
\newcommand{\notebox}[1]{\callout{notebg}{Note}{#1}}
\newcommand{\warnbox}[1]{\callout{warnbg}{Important}{#1}}

\title{\vspace{-1.5cm}\textbf{Android stale-cache diagnosis}\\[4pt]\Large Inflammation Tracker PWA\\[2pt]\normalsize Notes for follow-up --- 2026-06-11}
\author{}
\date{}

\begin{document}
\maketitle
\thispagestyle{fancy}
\vspace{-2em}

{\small\color{gray}\hrule\vspace{4pt}Working notes covering why an Android device may continue serving the previous version of \texttt{index.html}, \texttt{app.js}, or \texttt{sync.js} after a code change is deployed. Captured for later because the symptom recurred on 2026-06-11 after the \texttt{naproxen} column was added.\vspace{4pt}\hrule}

\vspace{1em}

\tableofcontents

% ============================================================
\section{Symptom}
% ============================================================

After deploying updated app code (e.g.\ a new form field, a CSS tweak, a logic fix), the Android phone continues to render the previous version even after force-closing the PWA and re-opening it. Desktop browsers pick up the change on next reload; Android does not.

This is \emph{not} the same as the v11-era ``\texttt{.htaccess} basic-auth stranded SW'' trap --- the service worker itself is already mitigated for that case (tolerant install, network-first fetch, non-OK never poisons the cache). The current symptoms point at caching layers \emph{above} the service worker, on the device.

% ============================================================
\section{Three suspects}
% ============================================================

\subsection{Suspect 1 --- HTTP cache underneath the service worker}

The service worker fetch handler in \texttt{app/sw.js} does:

\begin{lstlisting}[language=js]
const res = await fetch(req);
\end{lstlisting}

No \texttt{cache} option is supplied. The browser's HTTP cache sits \emph{below} the service worker: when the SW does \texttt{fetch(req)}, that request is allowed to be satisfied from the disk cache before ever touching the network. So ``network-first'' from the SW's point of view can in practice be ``HTTP-cache-first''. Android Chrome is more aggressive about this than desktop browsers, especially for assets with stable URLs and no cache-busting query string (\texttt{app.js}, \texttt{sync.js}, \texttt{index.html} all have stable URLs in this repo).

\textbf{How to check:} Attach the Android phone to a Mac with a USB cable and open \texttt{chrome://inspect} in desktop Chrome. With the PWA open on the phone, click \emph{inspect} to attach DevTools. Then:
\begin{itemize}[leftmargin=1.6em]
  \item DevTools $\to$ Network: reload the page and look at the \emph{Size} column for \texttt{app.js}, \texttt{sync.js}, \texttt{index.html}. If the column shows \texttt{(disk cache)} or \texttt{(memory cache)}, the SW handed over a cached copy without re-validating with the server.
  \item DevTools $\to$ Application $\to$ Cache Storage: expand \texttt{inflam-shell-vXX}. If \texttt{vXX} is older than the latest \texttt{CACHE\_VERSION} in \texttt{app/sw.js}, the new SW has not yet activated --- jump to Suspect 2.
  \item DevTools $\to$ Application $\to$ Service Workers: ensure status is \emph{activated and is running}. If a worker is listed as \emph{waiting}, the new SW is installed but the old one still controls the page; close all tabs/instances of the PWA, then re-open.
\end{itemize}

\textbf{Proposed fix (one line):} change the fetch handler in \texttt{app/sw.js} to bypass the HTTP cache:

\begin{lstlisting}[language=js]
const res = await fetch(req, { cache: "no-store" });
\end{lstlisting}

This forces every shell-asset fetch to go to the network, not the HTTP cache. The SW's own cache (the \texttt{Cache Storage} entry it maintains as an offline fallback) is unaffected and still wins when the network is unreachable or returns a non-OK status. The trade-off is one extra round-trip per asset per load --- negligible for seven small files.

\subsection{Suspect 2 --- the service worker script itself is cached}

The service worker is registered in \texttt{app/app.js}:

\begin{lstlisting}[language=js]
navigator.serviceWorker.register("sw.js");
\end{lstlisting}

The default value of \texttt{updateViaCache} is \texttt{"imports"}, which means the browser is allowed to fetch \texttt{sw.js} \emph{itself} through the HTTP cache. If the device's HTTP cache is holding a stale copy of \texttt{sw.js}, the browser never sees the new \texttt{CACHE\_VERSION} value, so the new SW never installs, never activates, and never wipes the old cache.

\textbf{How to check:} Same DevTools session as above, Application $\to$ Service Workers. The \emph{Source} link next to the current SW will show its contents. If \texttt{CACHE\_VERSION} in that source is older than the value currently in the repo, the SW script itself is stale.

\textbf{Proposed fix (one line):}

\begin{lstlisting}[language=js]
navigator.serviceWorker.register("sw.js", { updateViaCache: "none" });
\end{lstlisting}

With \texttt{updateViaCache: "none"}, the browser always bypasses the HTTP cache when checking for an updated \texttt{sw.js}. This is the canonical fix for ``my new service worker is never activating''.

\subsection{Suspect 3 --- server-side \texttt{Cache-Control} headers}

If the \texttt{.htaccess} file on the Hostinger deployment is setting a long \texttt{max-age} for \texttt{*.js} or \texttt{*.html} (or an \texttt{ExpiresByType} directive with a long default), then the HTTP cache is being \emph{told} by the server to hold these files for hours or days. This is the upstream cause of Suspect 1: even with \texttt{cache: "no-store"} in the SW, if a long-lived response is already in cache from before the fix, it will not be re-fetched until the cache entry naturally expires.

\textbf{How to check:} Two options.
\begin{itemize}[leftmargin=1.6em]
  \item On the server, view \texttt{.htaccess} and grep for \texttt{Cache-Control}, \texttt{ExpiresActive}, \texttt{ExpiresByType}, or \texttt{Header set Cache-Control}.
  \item From any browser, open DevTools $\to$ Network, reload the page, click \texttt{app.js}, look at the \emph{Response Headers} pane. Note the \texttt{Cache-Control} and \texttt{Expires} values that the server is sending.
\end{itemize}

\textbf{Proposed fix} (server side --- only if Suspect 3 turns out to be present): add the following to the \texttt{.htaccess}, scoped to the app's shell files:

\begin{lstlisting}[language={}]
<FilesMatch "\.(html|js|css|json)$">
  Header set Cache-Control "no-cache, must-revalidate"
</FilesMatch>
\end{lstlisting}

\texttt{no-cache} (despite the name) does not forbid caching --- it forbids \emph{using} a cached copy without revalidating with the server first. Combined with the SW's \texttt{cache: "no-store"} fetch, this guarantees the device sees the latest deployed code on every page load.

% ============================================================
\section{Recommended order of work}
% ============================================================

\begin{enumerate}[leftmargin=1.6em]
  \item Apply Suspect 1 and Suspect 2 fixes together --- both are one-line client-side changes in this repo, low risk, and address the most common Android-stale-PWA failure modes. Bump \texttt{CACHE\_VERSION} once after the change so existing devices flush their old SW-managed cache on the next visit.
  \item Re-test on the Android phone using the DevTools inspection steps above. In particular, watch for \texttt{app.js} loading with a non-cache \emph{Size} value, and \texttt{CACHE\_VERSION} in the running SW matching the repo.
  \item Only if the symptom persists, investigate Suspect 3 via the \texttt{.htaccess} on Hostinger.
\end{enumerate}

\warnbox{Be wary of one false positive when testing: if a PWA is installed to the home screen and was last opened with the old code, the splash transition can briefly flash the old UI from a screenshot before the new code paints. Wait a full second after the splash before judging whether the new code is live.}

% ============================================================
\section{Reference: where the relevant code lives}
% ============================================================

\begin{center}
\small
\begin{tabular}{lll}
\textbf{File} & \textbf{Function / Line} & \textbf{Purpose} \\ \hline
\texttt{app/sw.js}  & \texttt{fetch} listener (around L72--L104) & SW fetch handler, network-first \\
\texttt{app/sw.js}  & \texttt{CACHE\_VERSION} constant (L20)      & Cache namespace; bump to force eviction \\
\texttt{app/app.js} & \texttt{serviceWorker.register} (L479)      & SW registration call \\
\end{tabular}
\end{center}

\notebox{Line numbers are accurate as of 2026-06-11. They drift with edits --- search by name (\texttt{CACHE\_VERSION}, \texttt{serviceWorker.register}) if the lines no longer match.}

\end{document}
